\section{Atividades}

\begin{itemize}
      \item \textbf{Atividade 1:} Coleta, limpeza e pré-processamento dos dados, tratamento de inconsistências/valores ausentes e estruturação temporal dos registros da base fipeX por veículo.
      
      \item \textbf{Atividade 2:} Análise exploratória e pareamento comparativo visualização das trajetórias de depreciação, cálculo de estatísticas descritivas e pareamento dos modelos elétricos com equivalentes a combustão para identificação de padrões preliminares.
      
      \item \textbf{Atividade 3:} Transformação e representação espacial (\textit{reshaping}), realizando a codificação das séries temporais 1D em representações matriciais 2D utilizando as técnicas de GAF e MTF.
      
      \item \textbf{Atividade 4:} Modelagem de agrupamento e extração de padrões aplicando algoritmos de \textit{clustering} (K-Means e abordagens baseadas em densidade) sobre as representações geradas por imagem e sobre os dados vetoriais clássicos. Para garantir uma validação robusta da abordagem bidimensional proposta, o projeto estabelecerá uma linha de base (\textit{baseline}) operando diretamente sobre as séries temporais unidimensionais (1D) da taxa de retenção de valor. Serão aplicadas medidas de similaridade clássicas, com ênfase na Distância Euclidiana e, principalmente, no algoritmo de \textit{Dynamic Time Warping} (DTW) \cite{sakoe1978dynamic}, que é especialista em alinhar defasagens temporais. O desempenho dos agrupamentos obtidos no espaço 1D com DTW servirá como parâmetro de controle crítico para avaliar o real ganho de informação e precisão proporcionado pelas codificações matriciais 2D (GAF e MTF).

      \item \textbf{Atividade 5:} Avaliação experimental, interpretação e escrita final com o cálculo das métricas de validação de agrupamento, análise dos arquétipos comportamentais descobertos e consolidação dos resultados no artigo final.
\end{itemize}
